\documentclass[11pt,a4paper]{article}
\usepackage[margin=1in]{geometry}
\usepackage[T1]{fontenc}
\usepackage{lmodern}
\usepackage{amsmath,amssymb,booktabs,graphicx}
\usepackage[colorlinks=true,urlcolor=blue,citecolor=blue,linkcolor=blue]{hyperref}
\title{Finite-Size Scaling of Entanglement Across a Bond Defect\\
in the Critical Transverse-Field Ising Chain}
\author{Hugo Wei}
\date{Working research note --- October 2026 (preliminary, not peer reviewed)}

\begin{document}
\maketitle

\begin{abstract}
We investigate ground-state bipartite entanglement in a finite open
transverse-field Ising chain at criticality with a tunable central bond.
The model admits an exact free-fermion representation, allowing entropy
to be computed from a restricted Majorana covariance matrix rather
than from the exponentially large spin Hilbert space. We construct a
reproducible numerical baseline, independently validate the method
against exact diagonalization for small systems, and estimate the
coefficient of logarithmic entanglement scaling across defects.
The continuously varying asymptotic coefficient is a previously
established result; our ongoing investigation concerns finite-size
bias, fit-window dependence, and reliable numerical extrapolation.
This note is a preliminary research reproduction, not a claim of a
new universal law.
\end{abstract}

\section{Introduction}
Entanglement is a probe of the collective ground-state structure of
many-body quantum systems. At one-dimensional quantum critical points,
conformal field theory (CFT) constrains its logarithmic scaling
\cite{CalabreseCardy2004}. A local defect offers a tractable
perturbation: how much of the critical entanglement survives across
a bond whose strength varies between zero and the homogeneous limit?
The critical Ising defect problem is exactly understood in the
thermodynamic limit \cite{EislerPeschel2010}. Our immediate goal is a
careful computational reproduction and an assessment of finite-size
systematic errors.

\section{Model and theoretical background}
For even chain length $L$ with open boundary conditions, consider
\begin{equation}
 H=-\sum_{j=1}^{L-1}J_j\sigma_j^x\sigma_{j+1}^x
 -h\sum_{j=1}^L\sigma_j^z,\qquad
 J_j=\begin{cases}tJ,&j=L/2,\\ J,&\text{otherwise}.\end{cases}
 \label{eq:H}
\end{equation}
We set $J=h=1$, the bulk Ising critical point, and use
$0\leq t\leq1$ for the central bond. The subsystem $A$ comprises
the first $L/2$ spins; its von Neumann entropy is
\begin{equation}
S_A=-\mathrm{Tr}\,\rho_A\ln\rho_A,\qquad
\rho_A=\mathrm{Tr}_{\bar A}|\Psi_0\rangle\langle\Psi_0|.
\end{equation}
For the clean chain ($t=1$), boundary CFT predicts
\begin{equation}
S_A=\frac{c}{6}\ln L+b+o(1),\qquad c=\frac12.
\label{eq:CFT}
\end{equation}
The defect-dependent leading logarithmic coefficient can be
parameterized as $c_{\mathrm{eff}}(t)/6$. It is important not to
interpret a finite-$L$ fitted coefficient as a changed bulk central
charge: $c=1/2$ remains the Ising bulk value.

\section{Numerical method}
The Jordan--Wigner transformation maps Eq.~\eqref{eq:H} to free
Majorana fermions $\gamma_a$,
\begin{equation}
H=\frac{i}{4}\sum_{a,b=1}^{2L}A_{ab}\gamma_a\gamma_b,
\end{equation}
where $A$ is real and antisymmetric. In the convention used by the
code, $A_{2j-1,2j}=2h$ and
$A_{2j,2j+1}=2J_j$, with opposite antisymmetric entries.
For a finite chain without zero modes the ground-state covariance is
\begin{equation}
\Gamma_{ab}=\frac{i}{2}\langle[\gamma_a,\gamma_b]\rangle,
\qquad \Gamma=-A(-A^2)^{-1/2}.
\end{equation}
Let $\{\pm\nu_k\}_{k=1}^{\ell}$ be eigenvalues of
$i\Gamma_A$, where $\Gamma_A$ is the $2\ell$-dimensional
covariance restricted to the first $\ell$ sites. Then
\begin{equation}
S_A=\sum_{k=1}^\ell
\left[-p_k\ln p_k-(1-p_k)\ln(1-p_k)\right],
\quad p_k=\frac{1+\nu_k}{2},\quad \nu_k\geq0.
\end{equation}
The method scales polynomially with $L$. It avoids direct
diagonalization of the $2^L$-dimensional spin Hamiltonian.

\section{Validation and preliminary results}
We test antisymmetry of $A$, the pure-state identity
$\Gamma^2=-I$, vanishing cross-cut entropy when $t=0$, agreement
with independent exact diagonalization at $L=6$, and the expected
clean critical coefficient. These checks are implemented in
\texttt{tests/test\_ising.py}.

A pilot fit over $L=16,24,32,48,64$ uses
$S(L/2)=\frac{c_{\mathrm{fit}}(t)}6\ln L+b_t$:
\begin{center}
\begin{tabular}{ccc}
\toprule
Defect $t$ & $c_{\mathrm{fit}}(t)$ & Interpretation\\
\midrule
0 & $0.0000$ & halves disconnected\\
0.25 & $0.1532$ & finite-size estimate\\
0.50 & $0.3609$ & finite-size estimate\\
0.75 & $0.4881$ & finite-size estimate\\
1.00 & $0.5294$ & compare with $c=0.5$\\
\bottomrule
\end{tabular}
\end{center}
The excess at $t=1$ shows that short chains and simple
log-linear fits retain systematic corrections. These numbers
are \emph{not} estimates with extrapolated uncertainties and
should not be reported as asymptotic values.

\IfFileExists{../figures/entropy_vs_logL.png}{
\begin{figure}[ht]
\centering
\includegraphics[width=.85\linewidth]{../figures/entropy_vs_logL.png}
\caption{Finite-size entropy curves generated by
\texttt{python scripts/run\_experiment.py}.}
\end{figure}
}{
\noindent\fbox{\parbox{.94\linewidth}{Figure pending:
run \texttt{python scripts/run\_experiment.py} and compile again.}}
}

\section{Planned extension and limitations}
The next step is a quantitative comparison with the exact
defect-dependent expression derived by Eisler and Peschel
\cite{EislerPeschel2010}. We will vary fit windows, chain lengths,
and correction-to-scaling ans\"atze, assessing their stability.
A credible conclusion will require uncertainty analysis and
verification of the analytical formula's normalization and defect
parameter convention. Polynomial-time numerics do not by themselves
establish novelty.

\section{Conclusion}
The free-Majorana method provides a reproducible, independently
tested starting point for studying entanglement scaling in a
critical Ising chain with a single defect. The known qualitative
defect dependence is readily recovered, but finite-size fitting
produces systematic bias even in the homogeneous limit.
Further results will be added only after comparison with exact
literature and convergence analysis.

\paragraph{Status and reproducibility.}
This is a preliminary research note, prepared for independent
study and revision. Source code, tests, and numerical scripts
are maintained alongside the manuscript.

\bibliographystyle{unsrt}
\bibliography{references}
\end{document}
